%!TEX root = ../main.tex
% =============================================================
%  CAPÍTULO 7 — PRESENTACIÓN DE RESULTADOS
%                                       (Entrega 7 · 5 % · 13/10)
% =============================================================

\chapter{Presentación de resultados}\label{cap:resultados}

\begin{guia}[Guía de la cátedra: qué se evalúa en este capítulo]
\begin{enumerate}
    \item \textbf{Requerimientos formales}: \gls{rf} y \gls{rnf} con ID,
          descripción, prioridad y criterio de aceptación medible,
          trazados a los indicadores del capítulo~\ref{cap:metodologia}.
    \item \textbf{Evidencia de implementación}: capturas reales del sistema
          funcionando, enlaces al repositorio (anexo~\ref{anx:repositorios})
          y al despliegue.
    \item \textbf{Resultados cuantitativos honestos}: diseño experimental,
          métricas, matrices de confusión o su equivalente en el dominio,
          incluyendo lo que salió mal y su análisis. Reportar también lo
          malo y argumentarlo vale más que esconderlo. Hiperparámetros y
          versiones en el apéndice~\ref{apx:tecnicos}.
    \item \textbf{Validación con usuarios y expertos}: idealmente
          triangulada (métricas técnicas, encuesta, entrevista) y
          traducida a decisiones de diseño concretas. Instrumentos en los
          apéndices \ref{apx:encuesta} y \ref{apx:entrevista}.
\end{enumerate}
El sistema tiene que existir y andar: la entrega final incluye el
proyecto funcional.
\end{guia}

Este capítulo presenta los resultados obtenidos hasta esta entrega sobre
seis ejes: el modelo predictivo de severidad, el clustering de zonas de
riesgo, la validación con usuarios reales, la suite de pruebas
automatizadas, la ingesta de noticias de prensa y el comparador de
rutas. La hipótesis de trabajo del proyecto (1.5) sostiene que la
utilidad de esta herramienta no está en predecir con precisión un
siniestro individual, sino en la trazabilidad de las fuentes que
integra; los resultados de este capítulo se leen en función de qué tanto
la sostienen. El clustering (7.2) aporta la base de esa lectura:
confirma que el riesgo vial tiene patrones geográficos reales y
estables, sobre los que esa trazabilidad tiene algo concreto que
mostrar; el resto de los ejes puntúan si declarar procedencia,
antigüedad y cobertura resulta, en la práctica, más útil que un único
indicador de riesgo.

\section{Resultados del modelo predictivo de severidad}

El modelo vigente es un Random Forest de 300 árboles entrenado sobre
23.338 siniestros de moto georreferenciados en CABA, que estima la
gravedad probable de un siniestro (leve, grave o mortal) en una zona
dada, condicionada a que ese siniestro ocurra; no estima la probabilidad
de que un siniestro ocurra. Midió dos cifras de desempeño bajo dos
protocolos distintos, y la diferencia entre ambas es en sí misma un
resultado relevante del proyecto.

\begin{longtable}[]{@{}
  >{\raggedright\arraybackslash}p{(\columnwidth - 4\tabcolsep) * \real{0.3545}}
  >{\raggedright\arraybackslash}p{(\columnwidth - 4\tabcolsep) * \real{0.3213}}
  >{\raggedright\arraybackslash}p{(\columnwidth - 4\tabcolsep) * \real{0.3242}}@{}}
\caption{F1 macro del modelo de severidad según protocolo de evaluación.}\label{tab:7-1}\\
\toprule\noalign{}
\begin{minipage}[b]{\linewidth}\raggedright
\textbf{Protocolo de evaluación}
\end{minipage} & \begin{minipage}[b]{\linewidth}\raggedright
\textbf{F1 macro}
\end{minipage} & \begin{minipage}[b]{\linewidth}\raggedright
\textbf{Observación}
\end{minipage} \\
\midrule\noalign{}
\endfirsthead
\toprule\noalign{}
\begin{minipage}[b]{\linewidth}\raggedright
\textbf{Protocolo de evaluación}
\end{minipage} & \begin{minipage}[b]{\linewidth}\raggedright
\textbf{F1 macro}
\end{minipage} & \begin{minipage}[b]{\linewidth}\raggedright
\textbf{Observación}
\end{minipage} \\
\midrule\noalign{}
\endhead
\bottomrule\noalign{}
\endlastfoot
Split aleatorio interno & 0,368 & Mezcla años; usado durante la
selección de configuraciones \\
Evaluación temporal retrospectiva & 0,3032 & Año evaluado completo
(2024), nunca visto en entrenamiento. Corresponde a la variante sin edad
(con sexo e índice de inseguridad) vigente en producción \\
\end{longtable}

La caída de 0,368 a 0,3032 confirma que el primer número estaba inflado
por fuga de información entre años: el split aleatorio deja datos del
mismo período repartidos entre entrenamiento y evaluación, así que el
modelo termina memorizando patrones propios de ese período en vez de
generalizar. La evaluación temporal, que entrena con los años anteriores
y mide contra un año completo que el modelo nunca vio, da la cifra que
de verdad importa para decidir si el sistema sirve fuera del período con
el que se entrenó.

La matriz de confusión de esa evaluación temporal muestra dónde está el
límite real del modelo: en las clases más graves, la precisión es baja
(GRAVE 0,084; MORTAL 0,072) mientras que el recall es más alto (0,600 y
0,512 respectivamente). El modelo tiende a marcar muchas zonas como
graves o mortales de las que en realidad no lo son tanto, antes que
dejar pasar una zona grave sin marcar. Para un sistema de alerta
preventiva ese sesgo es defendible (es preferible sobre-advertir antes
que ignorar un caso grave), pero no habilita usar el score como una
medida fina de probabilidad: el sistema lo usa solo como criterio de
priorización relativa entre zonas, nunca como una cifra exacta de
riesgo.

Las ablaciones corridas sobre la variante anterior del modelo (que
todavía incluía la edad de la víctima) habían mostrado que sacarla como
variable subía el F1 macro de evaluación temporal a 0,286, y sacarla
junto con el sexo lo subía a 0,362, ambas por encima del 0,262 con
demografía completa; la demografía disponible (45\% de edad desconocida
en entrenamiento, 12\% en el año evaluado) no generalizaba bien entre
años. A partir de ese hallazgo, el modelo se reentrenó sin la variable
edad (conservando sexo e índice de inseguridad), vigente en producción;
los números de la Tabla 7.1 ya corresponden a esa variante. Esto es
consistente con un problema conocido en la predicción de severidad de
siniestros viales, en general difícil de resolver con buena precisión en
las clases más graves porque son, por definición, las menos frecuentes
en cualquier dataset real.

\section{Resultados del clustering de zonas de riesgo}

El universo de zonas de riesgo pasó por dos metodologías de clustering
antes de llegar a la que está en producción. El primer intento con
DBSCAN (eps=300m, min\_samples=8) fusionó casi toda la ciudad en un solo
cluster por efecto de encadenamiento, típico en tramas urbanas densas.
Ajustar el radio a 60m mejoró la separación (430 hotspots) pero dejó un
67,9\% de los siniestros sin cluster asignado.

La migración de DBSCAN a HDBSCAN, que no exige un radio fijo, es el
resultado que de verdad resolvió el problema: 412 a 740 zonas de
siniestros de moto, con la cobertura subiendo de 49,9\% a 75,6\% de los
datos. El mismo método aplicado a inseguridad general, con los
parámetros definitivos del sweep (mcs=150, ms=20), dio 633 hotspots con
50,2\% de cobertura. La brecha de cobertura entre ambas capas (75,6\%
contra 50,2\%) no es un error de calibración: responde a que los delitos
del GCBA están más dispersos geográficamente que los siniestros de moto,
que tienden a repetirse en las mismas intersecciones y avenidas.

\section{Resultados de la validación con usuarios reales}

Se aplicó un instrumento de validación a 12 motociclistas reales, en
sesiones individuales moderadas de forma virtual, con tareas de uso del
mapa y preguntas de percepción sobre una escala de 1 a 5. El
reclutamiento fue por contacto directo, así que los resultados son
indicativos y no representan a toda la población motociclista de
CABA/GBA. Los resultados centrales, para este capítulo, son tres.
Primero, el mapa en sí se usa sin dificultad: las tareas de activar una
capa y filtrar prensa se completaron sin ayuda en 10 de 12 casos.
Segundo, interpretar lo que el mapa muestra es harina de otro costal: la
tarea de comparar rutas y explicar la cobertura, la que exige entender
qué mide el número y no solo leerlo, solo la completaron sin asistencia
4 de 12, y 2 de 12 siguieron entendiendo el modelo como una predicción
de accidente pese a la aclaración textual existente. Tercero, la demanda
más repetida entre perfiles de uso muy distintos (reparto diario, uso de
fin de semana) fue la misma: reducir lo que se muestra por default y
dejar el detalle del cálculo como algo opcional.

Esa brecha entre uso del mapa y comprensión del score es el resultado
más accionable de esta validación: confirma que el riesgo del proyecto
no está en la interfaz, sino en la distancia entre lo que el sistema
declara medir y lo que una persona sin entrenamiento previo entiende al
usarlo. La muestra es chica y no aleatoria (n=12), así que estos
resultados son indicativos, no una medición poblacional.

\section{Resultados de la suite de pruebas automatizadas}

Al cierre de esta entrega, la suite de pruebas automatizadas tiene 168
pruebas de backend y 111 de frontend, 279 en total, todas en verde,
incluida la corrida en integración continua contra una base
PostgreSQL/PostGIS real y desechable (detalle en 4.10). La suite cubre
lógica de negocio, endpoints, cálculo de score de rutas y clasificación
de riesgo; no cubre todavía pruebas end-to-end contra navegador real ni
contra los servicios externos en producción (TomTom, clima), que siguen
siendo una verificación manual.

\section{Resultados de la ingesta de prensa}

Dos corridas reales y comparables muestran el efecto de un cambio de
fuente sobre el volumen de datos.

\begin{longtable}[]{@{}
  >{\raggedright\arraybackslash}p{(\columnwidth - 8\tabcolsep) * \real{0.2659}}
  >{\raggedright\arraybackslash}p{(\columnwidth - 8\tabcolsep) * \real{0.1773}}
  >{\raggedright\arraybackslash}p{(\columnwidth - 8\tabcolsep) * \real{0.1883}}
  >{\raggedright\arraybackslash}p{(\columnwidth - 8\tabcolsep) * \real{0.1883}}
  >{\raggedright\arraybackslash}p{(\columnwidth - 8\tabcolsep) * \real{0.1801}}@{}}
\caption{Candidatas, ubicadas, guardadas y relacionadas con moto por corrida del scraper.}\label{tab:7-2}\\
\toprule\noalign{}
\begin{minipage}[b]{\linewidth}\raggedright
\textbf{Corrida}
\end{minipage} & \begin{minipage}[b]{\linewidth}\raggedright
\textbf{Candidatas}
\end{minipage} & \begin{minipage}[b]{\linewidth}\raggedright
\textbf{Ubicadas}
\end{minipage} & \begin{minipage}[b]{\linewidth}\raggedright
\textbf{Guardadas}
\end{minipage} & \begin{minipage}[b]{\linewidth}\raggedright
\textbf{De moto}
\end{minipage} \\
\midrule\noalign{}
\endfirsthead
\toprule\noalign{}
\begin{minipage}[b]{\linewidth}\raggedright
\textbf{Corrida}
\end{minipage} & \begin{minipage}[b]{\linewidth}\raggedright
\textbf{Candidatas}
\end{minipage} & \begin{minipage}[b]{\linewidth}\raggedright
\textbf{Ubicadas}
\end{minipage} & \begin{minipage}[b]{\linewidth}\raggedright
\textbf{Guardadas}
\end{minipage} & \begin{minipage}[b]{\linewidth}\raggedright
\textbf{De moto}
\end{minipage} \\
\midrule\noalign{}
\endhead
\bottomrule\noalign{}
\endlastfoot
Histórica, con agregador (incluía Google Noticias) & 560 & 289 & 267 &
99 \\
Directa, tras retirar Google Noticias como fuente & 122 & 57 & 4 & 20 \\
\end{longtable}

La caída de 267 a 4 notas guardadas no es una regresión del sistema: es
la consecuencia directa y esperada de retirar Google Noticias como
fuente al confirmar que esa fuente no autorizaba ese uso. Esa fuente
aportaba la mayor parte del volumen de la corrida histórica, así que
perderla bajó el número bruto de notas, pero corregir una fuente no
autorizada era la decisión correcta aunque empeorara la métrica
superficial. Esa comparación no es un experimento controlado: la segunda
corrida parte de una base que ya tenía buena parte de esas notas
cargadas de corridas anteriores, así que contrasta dos corridas en
estados de base distintos, no un antes/después aislado del cambio de
fuente. Corridas posteriores con medios locales verificados (sitemaps
propios, sin agregador) vienen recuperando volumen de forma más lenta
pero legítima.

\section{Resultados del comparador de rutas}

El comparador de rutas calcula, para cada alternativa que devuelve el
buscador de caminos en tiempo real, un puntaje de riesgo a partir de
cuántas zonas de siniestros e inseguridad atraviesa el trazado y qué tan
cerca pasa de cada una. El puntaje no se usa en términos absolutos sino
relativos: las alternativas de un mismo pedido se comparan entre sí y se
etiquetan como bajo, medio, alto o crítico según su posición dentro de
ese grupo, no según un valor fijo, así que una ruta no cambia de
categoría solo porque el cálculo sumó o restó puntos por un ajuste
interno, mientras siga siendo, en términos relativos, la más o la menos
riesgosa del lote.

Se armaron cuatro casos reales contra el sistema en funcionamiento,
elegidos a partir de los propios datos de siniestros e inseguridad y no
al azar. Un caso de desvío, donde la ruta más corta de las tres
alternativas atraviesa una zona de inseguridad clasificada como crítica
y el sistema recomienda en cambio una opción 4,2 km más larga con
exposición baja. Un caso de cobertura parcial, donde una de las
alternativas cae fuera de la Ciudad de Buenos Aires y la cobertura de
inseguridad de esa zona baja a 56\% de los datos disponibles: un 44\%
del trayecto queda sin información de inseguridad para evaluarlo, no
porque no haya inseguridad ahí sino porque falta dato registrado. Un
caso de empate, donde tres alternativas llegan en el mismo tiempo
estimado con tráfico y el sistema las distingue solo por su nivel de
exposición al riesgo. Y un caso de discordancia entre fuentes, en una
zona con 2151 delitos registrados a diez metros del trazado pero apenas
12 siniestros de moto en el mismo tramo, donde las dos fuentes de datos
no coinciden en cuán riesgosa es la zona.

Sobre esos mismos cuatro casos se midió además qué tan sensible es el
resultado a dos parámetros internos del cálculo. El primero es el peso
relativo que se le da a cada nivel de riesgo al sumar el puntaje total
de una ruta: variarlo de una escala lineal a una que castiga mucho más
fuerte lo crítico no cambió en ningún caso cuál alternativa resultó la
de menor exposición (la recomendada), aunque el orden entre las
alternativas restantes sí puede cambiar según el peso elegido, porque la
comparación es por posición relativa dentro del lote, no por el valor
absoluto del puntaje. El segundo es la distancia máxima a la que una
zona de riesgo todavía se considera cruzada por la ruta: ahí sí hubo
cambios reales. Achicar esa distancia de 600 a 300 metros hizo caer la
cobertura de inseguridad del caso de cobertura parcial de 56\% a 18\%, y
en el caso de discordancia hizo que dos alternativas que antes se
distinguían con claridad (una de riesgo bajo, otra de riesgo medio)
terminaran con la misma clasificación. Agrandar esa misma distancia a un
kilómetro llegó a cambiar cuál de las tres alternativas del caso de
cobertura parcial salía mejor puntuada. Con estos resultados, el cálculo
es robusto al peso que se le da a cada nivel de riesgo, pero sensible al
umbral de distancia elegido, un parámetro que hasta esta medición no
tenía respaldo empírico propio más allá del criterio con que se fijó
originalmente.

La segunda medición de sensibilidad fue sobre el agrupamiento de tramos
de calle en corredores de riesgo, el proceso que decide qué calles
entran al mapa como corredor antes de mostrarlas. El parámetro que más
mueve el resultado es la cantidad mínima de siniestros que una calle
necesita para entrar en esa lista: bajarlo agregó un 50\% más de calles
candidatas y subirlo sacó más de la mitad de las que estaban, aunque el
orden de las calles que sí entran, de mayor a menor cantidad de
siniestros, no cambió en ninguna configuración. Un segundo parámetro, el
largo de tramo en que se divide cada calle para el análisis, no afecta
qué calles entran, solo en cuántos pedazos se las analiza: cambia el
nivel de detalle, no la cobertura.

Estos cuatro casos y las dos mediciones de sensibilidad quedan
documentados con datos reales, no estimados, y todavía están pendientes
de una revisión de criterio independiente sobre si cada resultado
concreto es razonable. Confirmar que el cálculo funciona como está
diseñado, que es lo que esta medición hace, es distinto de evaluar si
está diseñado de la forma más conveniente, que todavía no se hizo; esa
separación es deliberada en esta etapa del proyecto.

\section{Síntesis y vínculo con la hipótesis}

Los resultados de este capítulo sostienen la hipótesis del proyecto de
forma parcial, no absoluta. El clustering (7.2) sí encontró patrones
geográficos reales y estables, con una cobertura final razonable sobre
el dato disponible. El modelo (7.1) sí captura algo de señal, pero con
una capacidad de discriminación baja en las clases que más importan
(grave y mortal), así que su uso honesto es priorizar zonas entre sí, no
estimar una probabilidad precisa. La validación con usuarios (7.3)
agrega una dimensión que la hipótesis original no contemplaba: un
sistema puede capturar el patrón correctamente y aun así no ser
entendido por quien lo usa, y esa brecha de comprensión es, a esta
altura del proyecto, tan relevante como la precisión del modelo.

En conjunto, la evidencia reunida en este capítulo no permite afirmar
que MotoSafe GBA predice el riesgo vial de forma precisa, pero sí
sostiene que ordena el riesgo de forma consistente y verificable, con
sus límites declarados en vez de ocultos. Esa distinción, entre predecir
y priorizar con honestidad, es el resultado de fondo del proyecto hasta
esta entrega.
